\subsection{Свинки-копилки}

\subsubsection{Описание решения}

Чтобы определить минимальное число копилок, которое достаточно разбить, чтобы открыть остальные, введём структуру ориентированного графа \( G=(V,E) \) следующим образом. Вершины графа \( v\in V \) соответствуют копилкам, а рёбра вида \( (u,v)\in E \) эквивалентны связи «$u$ открывается ключом из $v$».

Возьмём произвольную вершину и покрасим её в некоторый цвет --- так мы отметим, что копилка была нами разбита. Далее найдём всех достижимых соседей вершины поиском в глубину, окрашивая обрабатываемые вершины в тот же выбранный цвет. Если мы изначально выбрали вершину \( M \) без входящих рёбер, то нам повезло, и мы закончили обработку данного подграфа с началом в этой вершине.

Но нам могло и не повезти, и на деле входящие рёбра у \( M \) всё же были. В таком случае при обработке очередного подграфа окажется, что мы пытаемся перекрасить \( M \) в другой цвет --- мы эту попытку пресекаем и не учитываем текущий подграф в аккумулируемой переменной, поскольку он является началом уже учтённого подграфа.

Если же в ориентированном графе есть неориентированные циклы, то мы смотрим на цвет окрашиваемой вершины. Если он совпадает с цветом текущего подграфа, то пропускаем эту вершину: эта ветка уже была обработана. 

Приведённый выше алгоритм нужно запустить для всех неокрашенных вершин в графе, чтобы в аккумулируемой переменной получить число наибольших ориентированных подграфов. Оно и будет соответствовать минимальному числу копилок.

\subsubsection{Асимптотическая сложность}

\textit{Временная сложность решения} --- линейная \( \mathcal{O}(N) \). Структура графа строится за линейное время относительно числа рёбер. Подсчёт ориентированных подграфов происходит при помощи итеративной реализации поиска в глубину, которая отрабатывает за линейное время относительно числа вершин и рёбер. Чтение и запись происходят за константное время.

\textit{Пространственная сложность решения} --- линейная \( \mathcal{O}(N) \). В памяти материализуется структура графа, которая хранит для каждой вершины константное количество данных. Суммарно число рёбер равно числу вершин, поэтому массивы соседних вершин не разрастаются.

\subsubsection{Листинг кода}

\begin{lstlisting}
#include <cstddef>
#include <iostream>
#include <list>
#include <unordered_map>
#include <vector>

struct Vertex {
  bool visited = false;
  std::size_t color = 0;
  std::vector<std::size_t> neighbours;
};

class Graph {
public:
  void AddEdge(std::size_t from_id, std::size_t to_id) {
    if (this->vertices_.count(from_id) == 0U) {
      this->vertices_[from_id] = {false, 0, {}};
    }
    if (from_id == to_id) {
      return;
    }
    this->vertices_[from_id].neighbours.push_back(to_id);
  }

  std::size_t CountSources() {
    std::size_t sources_count = 0;
    std::size_t current_color = 0;
    for (std::size_t i = 1; i <= this->vertices_.size(); i++) {
      if (this->vertices_[i].visited) {
        continue;
      }
      std::list<std::size_t> to_visit = {i};
      std::size_t broken_count = 1;
      while (!to_visit.empty()) {
        Vertex& top_vertex = this->vertices_[to_visit.front()];
        to_visit.pop_front();
        top_vertex.visited = true;
        top_vertex.color = current_color;
        for (const auto& neightbour : top_vertex.neighbours) {
          const Vertex& neightbour_vertex = this->vertices_[neightbour];
          if (neightbour_vertex.visited && neightbour_vertex.color != current_color) {
            broken_count = 0;
            continue;
          }
          if (neightbour_vertex.visited) {
            continue;
          }
          to_visit.push_front(neightbour);
        }
      }
      sources_count += broken_count;
      current_color++;
    }
    return sources_count;
  }

private:
  std::unordered_map<std::size_t, Vertex> vertices_;
};

int main() {
  std::size_t n_count = 0;
  std::cin >> n_count;
  Graph graph;
  for (std::size_t i = 1; i <= n_count; i++) {
    std::size_t open_idx = 0;
    std::cin >> open_idx;
    // inverting the direction
    graph.AddEdge(i, open_idx);
  }
  std::cout << graph.CountSources();
  return 0;
}
\end{lstlisting}